\part{Part I --- Light}

\chapter{What light requires of the medium}

Light is taken first because it is the sector whose linear content is most
nearly complete, and because its requirement on the medium is two-sided in a way
that constrains every later sector.

\section{S9: Light's Requirement on the Medium --- Shear on the Brane, None in
the Bulk}\label{step:S9}

\stagefield{Part / anchor}{Part I --- Light. Sector~1, step~1.}
\claimstatus{\StatusOpen}
\stagefield{Purpose}{State what light needs the medium to have. Light is a
transverse wave and a GNLS superfluid cannot carry one, so light's first act is
to name the structure the GNLS does not supply. The step introduces the two
constants the whole sector rests on and fixes the light cone as their ratio.}

\paragraph{Status.}
\StatusOpen. First step banked under the requirements-first direction: light
states what it needs, and the medium's structure is not assumed in advance. The
result is a \emph{requirement with a live falsifier}, not a derivation --- both
constants enter \textbf{postulated}, and whether one substructure can meet both
halves of the requirement is a knit question this step does not settle.

\paragraph{The argument, in three moves.}
\resultanchor{1 --- What Maxwell demands.} Two transverse polarisations, one
speed, no longitudinal mode.

\resultanchor{2 --- What a GNLS has.} Linearising Madelung about a uniform
\(\rho_0\) gives exactly \emph{one} mode,
\[
  \omega^2 = c_s^2k^2 + \frac{\hbar^2}{4m^2}k^4,
  \qquad
  c_s^2 = \frac{1}{m}\frac{dP}{d\rho} = \frac{5K\rho_0^4}{m},
\]
and it is \textbf{longitudinal}. The obstruction is not incidental but
representation-theoretic: \(\Psi=\sqrt{\rho}\,e^{i\theta}\) is one complex
scalar, its excitations \(\delta\rho\) and \(\delta\theta\) are both spin-0, and
one broken \(U(1)\) gives one Goldstone phonon --- so a scalar's spectrum cannot
contain a spin-1 photon, \emph{in any dimension}.

\resultanchor{3 --- Therefore.} Light requires structure the GNLS does not
contain. Brane elasticity is not a consequence of the GP/NLS mean field: a
GP/NLS superfluid is a \emph{fluid}, with zero shear modulus. The brane's shear
rigidity therefore requires the \textbf{substructure} beneath the mean field ---
the standing postulate that the GNLS is the effective, coarse-grained
description of a medium with constituents of its own.

\paragraph{Provenance of moves 2 and 3 --- one source, not two.}
Both quotations come from the same document,
\StageFile{software/stage1\_solver/decisions/15\_em\_medium\_native\_physical\_picture.md}
(lines 29--40 and 230--233). It is an external review, and \textbf{no script in
this repository executes either argument}. A weaker in-repo form exists ---
\(v=(\hbar/m)\nabla\theta \Rightarrow \nabla\times v\equiv 0\), so superfluid
flow is curl-free --- but it is never wired into a no-shear proof. The claim
that a GNLS cannot carry light therefore rests on a single external document
with no executing script; it is stated here at that strength and no higher.

\paragraph{The requirement --- and it is two-sided.}
\begin{center}
\begin{tabular}{@{}lll@{}}
\toprule
phase & requirement & what it buys \\
\midrule
ordered (\(\chi_B=1\))    & carries MacCullagh curl-only shear & photons \emph{exist} \\
disordered (\(\chi_B=0\)) & carries \emph{no} shear            & photons stay \emph{confined} \\
\bottomrule
\end{tabular}
\end{center}

\noindent
Both halves are light's own bill. The second is \emph{not} magnetism's: without
it a photon has somewhere else to go, light leaks off our three-dimensional
space, and that is an energy sink with no observational room. The requirement is
stressed hardest at the throat, where the brane is bent into \(\pm w\).

A \emph{third} consumer of the second half: the recorded throat-support
mechanism is an outward trapped standing-wave pressure whose trapped mode is a
brane-shear standing wave. If the bulk carried shear, that mode radiates away,
the outward pressure vanishes, and the throat closes. Bulk shear-freeness is
what makes the geon possible at all.

\paragraph{The equations.}
\[
  L=\tfrac12\rho_{br}(\partial_t u)^2-\tfrac12\mu_R(\nabla\times u)^2,
  \qquad
  c_\gamma^2=\frac{\mu_R}{\rho_{br}} \quad (R4).
\]

\paragraph{Dimensions, derived from the action.}
\(L\) is an energy density on the three-dimensional brane,
\([L]=M\,L^{-1}T^{-2}\), and \([u]=L\); hence \([\partial_t u]=LT^{-1}\) and
\([\nabla\times u]=1\), so
\[
  [\rho_{br}]=M\,L^{-3}=[-3,0,1],
  \qquad
  [\mu_R]=M\,L^{-1}T^{-2}=[-1,-2,1],
\]
and \(R4\) gives \([-1{-}({-}3),\,-2{-}0,\,1{-}1]=[2,-2,0]=2\times[1,-1,0]\), so
\(c_\gamma=[1,-1,0]\). These were derived from the Lagrangian \emph{before}
consulting the parameter register, which records the same two dimensions
independently (\StageFile{notes/parameter\_register.md}, lines 137--138).

\paragraph{What is new.}
\begin{center}
\begin{tabular}{@{}lll@{}}
\toprule
item & class & why \\
\midrule
\(\rho_{br}\) & \textbf{postulated} & a substructure property; not derivable from the GNLS \\
\(\mu_R\)     & \textbf{postulated} & ditto; the polar-\(P\) route returned \texttt{FAIL\_COUPLE\_STRESS\_NOGO} (B2) \\
\(c_\gamma\)  & \textbf{derived}    & \(R4\), from the two above \\
curl-only form & \textbf{postulated} & forced by Maxwell's no-longitudinal demand \\
\bottomrule
\end{tabular}
\end{center}

\noindent
The far field consumes only the ratio \(c_\gamma\). That must not be read as
``the two are not separately owed'': a \emph{simulation} needs \(\mu_R\) and
\(\rho_{br}\) absolutely.

\paragraph{Inputs, by locus.}
The standard GNLS (\(\psi\), \(\rho\), \(\theta\), Madelung,
\(v=(\hbar/m)\nabla\theta\)) comes standard with the vocabulary; Maxwell's three
demands come from the imported theory. \textbf{Nothing is taken from an earlier
v3 step} --- S9 is sector~1, step~1.

\paragraph{Regime.}
Linear response about an unstrained brane; small oscillations. Everything here
is quadratic.

\paragraph{Departure.}
None at this step. The departure enters at S11, when the medium's
compressibility overrides the no-longitudinal demand.

\paragraph{Falsifier, and its status.}
\emph{Requirement:} the substructure carries shear in its ordered phase and none
in its disordered phase. \emph{Falsifier:} a substructure supplying one and not
the other --- in particular one carrying shear in the \emph{disordered} phase,
which would unconfine light \emph{and} kill the trapped-mode geon.
\emph{Status:} \textbf{live}, and it is a knit question. Three independent
consumers (photon propagation, photon confinement, geon support), and nothing
yet shows one substructure can do both. Bulk shear-freeness is currently
\textbf{postulated}; the only machine check on it is dimensional.

This is light's first open requirement conflict, and it is \emph{internal to
light} --- not a light-versus-magnetism conflict. It is a question about the
ordering transition: can one substructure be ordered-and-shear-bearing in one
phase and unstructured-and-shear-free in the other?

\paragraph{Registry additions executed at this step.}
\(Q.brane.rho\_br\ [-3,0,1]\), \(Q.brane.mu\_R\ [-1,-2,1]\),
\(Q.brane.c\_gamma\ [1,-1,0]\), and relation \(R4\):
\(c_\gamma=\sqrt{\mu_R/\rho_{br}}\). \(R4\) is the corpus's existing name for
this relation. Note that \(c_\gamma\) re-enters \emph{here}, with provenance,
having been deleted from the \emph{medium} block at S0.5: that round trip is the
mechanism working.

% NOT \stagefield{Verification} -- that field is suppressed in the default
% build, and this is a GAP DISCLOSURE, not a verification listing.  It must
% remain visible.  Delete these two artifacts from the "owed" wording only when
% they actually exist.
\paragraph{Verification --- and what is missing.}
Registry gates only, at time of writing: dimensional homogeneity
\StatusText{pass}, acceptance \StatusText{match}, able-to-fail harness
\StatusText{pass}. \textbf{A step-level SymPy audit and an independent
Mathematica audit do not yet exist for v3.} \(R4\) is carried as a
\emph{declared} relation whose execution locus points into v2
(\StageFile{research/pde\_ledger\_v2/scripts/ledger\_stage003\_transverse\_photons\_stray\_longitudinal\_sympy\_audit.py},
lines 516--601). The gates check that \(R4\) is dimensionally homogeneous and
that its dataflow closes; \textbf{they do not derive
\(\omega^2=(\mu_R/\rho_{br})k^2\) from the Lagrangian.} Both audits are owed
before this step's six-artifact unit is complete.


\section{S10: Two Transverse Photons --- and the Count Comes from the
Brane's Own Dimension}\label{step:S10}

\stagefield{Part / anchor}{Part I --- Light. Sector~1, step~2.}
\claimstatus{\StatusText{Computed}}
\stagefield{Purpose}{Count the independent waves carried by the brane's
curl-only sector and identify why that count is two. The result is
\(D_{\mathrm{brane}}-1\) transverse modes, obtained from the nullity of the
dynamical matrix at each root. It is not inferred from a factorised
determinant.}

\paragraph{Status.}
\StatusText{Computed}. S9 supplied the curl-only cone; this step computes its
mode content. The count is banked, while the brane dimension that supplies the
number is explicitly \textbf{postulated}: \(D_{\mathrm{brane}}=3\) is what
light requires of the medium, not something derived here.

\paragraph{Identification, flagged before use.}
The field \(u\) is the \textbf{in-plane} displacement: a \(D\)-vector on the
brane. It is not a \(D+1\)-vector that includes motion out of the brane into
\(\pm w\). That out-of-plane motion is the separate \(h\)-branon that carries
charge, and \(h\ne u_L\). This identification is load-bearing: if out-of-plane
motion belonged to the same elastic sector, the transverse count at \(D=3\)
would be \textbf{three}, not two.

\paragraph{The argument, in three moves.}
\resultanchor{1 --- Decompose about \(k\).} For
\(u=a\exp[i(k\cdot x-\omega t)]\), the constant amplitude splits into one
component along \(\widehat{k}\) and a perpendicular component in a
\((D-1)\)-dimensional subspace:
\[
  D=1+(D-1).
\]

\resultanchor{2 --- The stiffness sees only the perpendicular part.} The
antisymmetric derivative sends \(\nabla\times u\) to \(i k\times a\), in the
three-dimensional notation, and in every \(D\) its quadratic form is
\[
  \frac12\sum_{i,j}(k_i a_j-k_j a_i)^2
    = |k|^2|a|^2-(k\cdot a)^2.
\]
It therefore penalises only \(a_\perp\). A longitudinal amplitude has zero
curl and zero energy cost. The longitudinal mode is \textbf{free, not
forbidden}: there is no restoring force because the medium resists twisting
and is indifferent to compression along the wave. An ordinary elastic solid,
which resists both, also carries a longitudinal wave.

\resultanchor{3 --- Solve per piece.} The equation of motion is
\[
  \rho_{\mathrm{br}}\omega^2 a
    =\mu_R\bigl[k^2a-k(k\cdot a)\bigr].
\]
It gives
\[
  \omega^2=0 \quad\hbox{for }a\parallel k,
  \qquad
  \omega^2=\frac{\mu_R}{\rho_{\mathrm{br}}}k^2
    \quad\hbox{for }a\perp k.
\]
The transverse modes are degenerate by \textbf{symmetry}, not coincidence:
nothing in the equation distinguishes one perpendicular direction from
another, so all carry the same frequency.

\paragraph{The computed count.}
The two audit engines evaluate the dynamical matrix at each root and compute
its \textbf{nullity}. The count below is therefore not read off the
multiplicity of a factor in a factorised determinant.
\begin{center}
\begin{tabular}{@{}cL{0.23\textwidth}L{0.39\textwidth}c@{}}
\toprule
\(D\) & \(\omega^2=0\) &
\(\omega^2=(\mu_R/\rho_{\mathrm{br}})k^2\) & sum \\
\midrule
2 & nullity 1, \(\parallel k\) & nullity 1, \(\perp k\) & 2 \\
\textbf{3} & nullity 1, \(\parallel k\) &
  \textbf{nullity 2}, \(\perp k\) & 3 \\
4 & nullity 1, \(\parallel k\) & nullity 3, \(\perp k\) & 4 \\
5 & nullity 1, \(\parallel k\) & nullity 4, \(\perp k\) & 5 \\
\bottomrule
\end{tabular}
\end{center}

\noindent
Thus the longitudinal nullity is one and the transverse nullity is \(D-1\).
In particular, exactly two transverse polarisations require
\(D_{\mathrm{brane}}=3\).

\paragraph{Why this is not a codimension count.}
The bulk never enters the computation. There is no codimension anywhere in
the dynamical matrix or its nullities; only the brane's own dimension appears.
Codimension is therefore not merely the wrong number, but an
\textbf{absent quantity}. The \(D=2,3,4,5\) pattern above is the evidence.

\paragraph{Equations and dimensions for general \(D\).}
The Lagrangian is an energy density on a \(D\)-dimensional brane,
\[
  L=\frac12\rho_{\mathrm{br}}(\partial_tu)^2
    -\frac12\mu_R(\operatorname{curl}u)^2,
  \qquad
  (\operatorname{curl}u)^2
    \equiv\frac12\sum_{i,j}(\partial_i u_j-\partial_j u_i)^2.
\]
The latter is the antisymmetric derivative in every \(D\). At \(D=3\), its
reduction to the ordinary \( |\nabla\times u|^2 \) is computed, with residual
zero, rather than assumed. The constituent dimensions close as functions of
\(D\):
\[
  [\rho_{\mathrm{br}}]=(-D,0,1),
  \qquad
  [\mu_R]=(2-D,-2,1).
\]
At \(D=3\) these become \((-3,0,1)\) and \((-1,-2,1)\), agreeing with the
registry.

These formulas explain an S9 loose end by identity. S9 found that changing the
assumed brane dimension corrupted both constituent dimensions while the speed
check stayed green. For every \(D\), however,
\[
  [\mu_R]-[\rho_{\mathrm{br}}]
    =(2-D)-(-D)=2
\]
in the length slot. The speed's dimension \emph{cannot} see the brane
dimension; the apparently weak check is a property of the physics, not an
unexplained audit blind spot.

\paragraph{What is new.}
\begin{center}
\begin{tabular}{@{}lll@{}}
\toprule
item & class & why \\
\midrule
\(D_{\mathrm{brane}}=3\) & \textbf{postulated} & what light requires; not derived \\
mode count & \textbf{derived} & nullity: \(D_{\mathrm{brane}}-1\) transverse, one longitudinal \\
\bottomrule
\end{tabular}
\end{center}

\noindent
The dimension closes a real hole: S9's dimensions depended on it, but it had
not been declared. The derived \(n_{\mathrm{pol}}\) is deliberately not a
separate registry row. The constraint machinery rejects relations whose
residual contains discrete symbols, and the closed schema has no vocabulary
for expressing that one discrete value derives from another. A bare row would
make a derived count look like a free choice, which is less faithful than
omitting it.

\paragraph{Inputs, by locus.}
The Lagrangian, \(\rho_{\mathrm{br}}\), \(\mu_R\), and the curl-only form come
from S9. The identification of \(u\) as the in-plane displacement is made and
flagged in this step. \textbf{Nothing comes from the bulk}; it does not enter.

\paragraph{Regime.}
Linear response about an unstrained brane; small oscillations.

\paragraph{Departure.}
None at this step. But \(\omega^2=0\) means the longitudinal slot is
\textbf{non-propagating}, not absent. S10 therefore does not deliver Maxwell's
no-longitudinal-photon demand; once the medium is allowed to be compressible,
that zero lifts. This is not a defect awaiting repair. Exact Maxwell would be
the failure: it puts charge in by hand, so a model matching it exactly would
have no way to anchor charge physically. The extra slot is that anchor and is
carried forward to S11. It must not be confused with the charge field itself:
\(u_L\) is in-plane, whereas the charge-bearing \(h\)-branon is out of plane
and \(h\ne u_L\).

\paragraph{Verification.}
Two independently written engines agree on every value. The blind Mathematica
audit
(\StageFile{research/pde\_ledger\_v3/mathematica/S10\_two\_transverse\_photons\_mathematica\_audit.wl})
reads nothing and was built first. The SymPy audit
(\StageFile{research/pde\_ledger\_v3/scripts/S10\_two\_transverse\_photons\_sympy\_audit.py})
imports the shared registry and was built while the Mathematica file was
quarantined out of the tree, so it could not transcribe from it. After
restoration, that file was byte-identical to its committed blob.

Every value matched predictions committed before either script ran
(\StageFile{research/pde\_ledger\_v3/steps/S10\_PREREGISTERED\_PREDICTION.md},
commit \texttt{bc276485}), including the two flagged uncertainties: the
antisymmetric form reduces to ordinary curl-squared at \(D=3\), and \(D=2\) is
not degenerate in kind.

The able-to-fail control changes the \emph{form}: replacing the
antisymmetric-derivative stiffness by full gradient-squared stiffness
collapses the two roots into one root of nullity \(D\), removes the
\(\omega^2=0\) root, and makes the longitudinal mode propagate. A coefficient
control cannot distinguish this step's claim: rescaling \(\mu_R\) moves the
roots but leaves their nullities unchanged. A coefficient control tests the
arithmetic; only the form control tests the physics of the count.

Registry gates were all green with the payload unchanged: acceptance
\StatusText{match} (\(10\mathbin{\to}6\), \(10\mathbin{\to}6\),
\(10\mathbin{\to}5\)), dimensional homogeneity \StatusText{pass},
able-to-fail \StatusText{pass}, and 11 tests.


\section{S11: The Stray Longitudinal --- and Every Question It Raises Is One
Question}\label{step:S11}

\stagefield{Part / anchor}{Part I --- Light. Sector~1, step~3.}
\claimstatus{\StatusText{Computed Spectrum / Open Physical Status}}
\stagefield{Purpose}{Lift S10's longitudinal zero when the brane is allowed to
resist compression, enter the brane compression modulus as a postulated knob
with a named retirement condition, and state what the resulting spectrum does
and does not establish.}

\paragraph{Status.}
The longitudinal zero is lifted and its cone is computed. Its physical status
is not: whether the mode radiates, is observable, or constitutes a
Lorentz-breaking cone remains open because the required brane--bulk interface
coupling law has not been built.

\paragraph{Departure --- one statement, not three.}
\textbf{The second mode's entire physical status --- radiative, observable,
Lorentz-breaking, or none of the above --- reduces to one unbuilt object: the
brane--bulk interface coupling law that says whether and how matter couples to
it.} The law is linear and was deferred by choice, not by difficulty; it is the
work of S11b. The three questions carried by the walk are therefore one:
\begin{center}
\begin{tabular}{@{}L{0.60\textwidth}L{0.27\textwidth}@{}}
\toprule
question & reduces to \\
\midrule
does the longitudinal radiate into the bulk, or stay bound? & the coupling law \\
is light's confinement unconditional, or does it rest on a
polarisation-overlap argument? & the coupling law \\
does a second characteristic speed break Lorentz invariance for us? & the coupling law \\
\bottomrule
\end{tabular}
\end{center}

\paragraph{What S11 does not deliver.}
S11 does \textbf{not} decide bound versus radiating; does not establish that
light's confinement is unconditional; does not establish whether the
longitudinal is observable or unobservable; and supplies no leakage rate.

\paragraph{Why the extra mode is not a defect being repaired.}
Exact Maxwell would be the \textbf{failure}: Maxwell puts charge in by hand, so
a model matching it exactly would have no way to anchor charge physically. The
extra mode is that anchor; it is what made the drum-head charge picture click
(\StageFile{research/pde\_ledger\_v3/CHARTER.md\#falsification-standard}).
The token \texttt{FAIL\_CAUCHY\_STRAY\_LONGITUDINAL} is misnamed; its prefix is
not a verdict.

\paragraph{Identification, flagged before use.}
The field \(u\) is the material displacement of the stuff whose density is
\(\rho_{\mathrm{br}}\). S9's kinetic term
\(\tfrac12\rho_{\mathrm{br}}(\partial_tu)^2\) is kinetic energy only if
\(\dot u\) is the velocity of the material carrying that inertia. Linearised
continuity therefore applies,
\[
  \delta\rho_{\mathrm{br}}=-\rho_{\mathrm{br}}\nabla\!\cdot u,
\]
and the medium's equation of state makes compression cost energy. This is not
a modelling choice. It would have been wrong if \(u\) were an
orientation/director field: then \(\nabla\!\cdot u\) would be splay rather than
a density change, and the channel would need a separate argument.

\paragraph{Correction 1 --- the invariant count and the method.}
The derivative \(\partial_i u_j\) decomposes into trace, symmetric-traceless,
and antisymmetric parts. S10 did not forget a term: it retained one invariant
of several. Curl-only
stiffness charges twist and nothing else, so the pure-trace longitudinal wave
was free rather than forbidden. But the walk's claim of ``exactly three
coefficients for every \(D\geq2\)'' was \textbf{wrong} under proper rotations:
\[
  N_{SO}=\{D=2\mapsto4,\ D=3\mapsto3,\ D=4\mapsto4,\ D=5\mapsto3\},
  \qquad N_O=3\quad\hbox{for every }D.
\]
The method, not its arithmetic, was the error. Decomposing into pieces that do
not mix and counting them misses cross-pairings between isomorphic summands.
At \(D=2\), the trace and \(\Lambda^2\) are both \(SO(2)\) scalars; at \(D=4\),
\(\Lambda^2\) splits into self-dual and anti-self-dual pieces. The additional
reflection-odd invariants are
\[
  (\operatorname{tr}G)\,\epsilon^{ij}G_{ij}\quad(D=2),
  \qquad
  \epsilon_{ijkl}G_{ij}G_{kl}\quad(D=4).
\]
Five independent derivations agree on the corrected counts. The \(D=2\)
invariant is not a total derivative: its Euler--Lagrange operator is non-zero,
so an \(SO(2)\)-invariant Lagrangian can mix longitudinal and transverse
sectors. The \(D=4\) invariant is a total derivative and adds no operator
structure. Thus sector separation is a \textbf{\(D=3\) fact}, not a structural
one.

% Every rendered occurrence carries the status label required by the record.
% The underbrace is an annotation, not part of the symbol or an extra index.
\newcommand{\SXIcomp}{\underbrace{B_{\mathrm{comp}}}_{\text{\scriptsize postulated}}}

\paragraph{The spectrum.}
With the brane compression modulus \(\SXIcomp\), postulated rather than
derived, the trace term enters strictly as
\[
  \rho_{\mathrm{br}}\omega^2a
  =\mu_R\bigl[k^2a-k(k\!\cdot\!a)\bigr]
   +\SXIcomp\,k(k\!\cdot\!a).
\]
For \(k\!\cdot\!a=0\), the new term vanishes identically: compression cannot
touch light, and S9/S10 are not reopened. For \(a\parallel k\), the
\(\mu_R\) bracket vanishes and the old zero is replaced, not perturbed.

The two independently written engines compute
\[
  M_{ij}=\mu_R(k\!\cdot\!k)\delta_{ij}
         +(\SXIcomp-\mu_R)k_i k_j,
\]
and obtain:
\begin{center}
\begin{tabular}{@{}L{0.24\textwidth}L{0.65\textwidth}@{}}
\toprule
quantity & computed value \\
\midrule
transverse &
\(\omega^2=(\mu_R/\rho_{\mathrm{br}})k^2\), nullity \(D-1\), kernel
perpendicular to \(k\); unchanged from S9/S10 \\
longitudinal &
\(\omega^2=(\SXIcomp/\rho_{\mathrm{br}})k^2\), nullity \(1\), kernel
parallel to \(k\) \\
cross-sector &
the perpendicular root is independent of \(\SXIcomp\); the parallel root is
independent of \(\mu_R\); both computed residuals are zero \\
degeneracy & exactly \(\SXIcomp=\mu_R\) \\
\bottomrule
\end{tabular}
\end{center}

The dimensions close in general \(D\) as
\[
  [\rho_{\mathrm{br}}]=(-D,0,1),
  \qquad
  [\mu_R]=[\SXIcomp]=(2-D,-2,1),
  \qquad
  [c_\gamma]=[c_L]=(1,-1,0).
\]
Thus the longitudinal speed is the derived quantity
\[
  c_L^2=\frac{\SXIcomp}{\rho_{\mathrm{br}}}.
\]
Every computed value agrees with the preregistration
(\StageFile{research/pde\_ledger\_v3/steps/S11\_PREREGISTERED\_PREDICTION.md},
commit \texttt{67d919bd}) except the move-2 invariant count: there the
preregistration was wrong and the engines were right.

\paragraph{Controls.}
The form control replaces the trace invariant with the
symmetric-traceless one and moves both roots:
\[
  \omega_\perp^2=\frac{\mu_R+\mu_{\mathrm{br}}}{\rho_{\mathrm{br}}}k^2,
  \qquad
  \omega_\parallel^2=\frac{4\mu_{\mathrm{br}}}{3\rho_{\mathrm{br}}}k^2.
\]
The transverse root thereby acquires \(\mu_{\mathrm{br}}\), showing that the
trace invariant is the unique one that lifts the longitudinal while leaving
light untouched. The coefficient control merely rescales \(\SXIcomp\) and
moves only the parallel root, so it cannot test that shape claim. The form
control's \(4/3\) independently reproduces the corpus's rejected-Cauchy-branch
coefficient from an engine blind to that document.

\paragraph{Correction to the compression mechanism.}
The walk's first mechanism was \textbf{wrong}: it had brane material
``squeezing out into the bulk,'' which smuggled in a phase transition because
the bulk is the disordered phase. Compression does not make the brane
unordered. The correct mechanism is that the brane \textbf{thickens}: in-plane
compression raises areal density, accommodated either as higher four-dimensional
density, charged by the equation of state, or as a wider wall, charged by the
wall potential. The order parameter never changes; the sheet bulges into
\(\pm w\). The two channels take additive shares of one strain, so they are
springs in series, their compliances add, and the postulated \(\SXIcomp\) is
softer than either channel alone.

\paragraph{The flat direction and the named retirement condition.}
S7 records that the double well selects no slab width. If that flatness
survived, \(B_{\mathrm{wall}}=0\), hence the postulated \(\SXIcomp=0\), and
S10's zero would not lift. Width gradients lift it: a wave modulates the width,
tilts and stretches the interfaces, and costs
\(\sigma_{\mathrm{wall}}|\nabla W|^2\). The mode is therefore flat at \(k=0\)
and stiff as \(k^2\), predicting a flexural crossover with
\(\omega\propto k^2\) at long wavelength.

The cone above was computed with wall width \textbf{frozen}; unfreezing it
should soften the mode. If it does not, move~5 was wrong and must not be
reconciled away. The wall tension \(\sigma_{\mathrm{wall}}\) is S6-derived, so
the named retirement condition is
\texttt{DEFECT\_REGISTER\#c12}: the postulated \(\SXIcomp\) is to be retired
visibly at S6 if that derivation supplies it rather than leaving an independent
knob. The present knob count is an upper bound that can only improve.

\paragraph{Correction 2 --- bulk matching was overstated in both directions.}
Phase matching gives
\[
  k_w^2=k^2\left(\frac{c_L^2}{c_{s0}^2}-1\right).
\]
The walk read \(c_L>c_{s0}\) as ``radiates'' and \(c_L<c_{s0}\) as ``bound.''
\textbf{Neither direction is established.} Phase matching is kinematic only:
it determines whether a propagating bulk channel exists, not whether the mode
uses it, and not whether an evanescent solution forms a bound eigenmode. Both
claims require the absent interface coupling law. The omitted third case is
\(k_w=0\), grazing.

\paragraph{What is new.}
\begin{center}
\begin{tabular}{@{}lll@{}}
\toprule
item & class & why \\
\midrule
\(Q.\mathrm{brane}.\SXIcomp\) & \textbf{postulated} & retire under
\texttt{DEFECT\_REGISTER\#c12}; knob count is an upper bound \\
\(Q.\mathrm{brane}.c_L\) & \textbf{derived} & \(R5\), from the postulated
\(\SXIcomp\) and \(\rho_{\mathrm{br}}\) \\
the mode itself & \textbf{derived} & nullity of the dynamical matrix \\
\bottomrule
\end{tabular}
\end{center}

\noindent
The postulated \(\SXIcomp\) is a \textbf{brane} compression modulus. It is not
\(K_{\mathrm{br}}\), the rejected elastic branch's bulk modulus; not
\(B_{\mathrm{eff}}=\rho_{B0}^2/\chi_c\); and not the bulk medium's modulus. It
is registered with \texttt{aliases: []} and no borrowed loci. The registry
changes are ambient \(10\mathbin{\to}12\) and residue \(6\mathbin{\to}7\), with
residue
\(\{\hbar,m,K,\rho_0,\rho_{\mathrm{br}},\mu_R,\SXIcomp\}\). The discrete
payload \(\{n_{\mathrm{eos}}=5,D_{\mathrm{brane}}=3\}\) is unchanged and
remains off the continuous axis.

\paragraph{Inputs, by locus.}
S9 supplies the material-displacement kinetic term and the postulated
\(\rho_{\mathrm{br}}\) and \(\mu_R\); S10 supplies the longitudinal zero and
the in-plane identification. The equation of state supplies the compression
cost. S7 supplies the flat slab-width statement, and the S6-derived wall
tension supplies the proposed retirement route. The interface law is not an
input because it has not been built.

\paragraph{Regime.}
Linear response about the brane, with wall width frozen for the computed cone.
The predicted long-wavelength softening belongs to the width-unfrozen problem
and is not computed here.

\paragraph{What this settles about the light sector as a whole.}
First, \(c_L\) is \textbf{not} a gravitational wave. Gravity in this model is
the flow between draining defects, carried by flow plus Bernoulli pressure and
not by ripples or radiation
(\StageFile{docs/conceptual\_foundation.md}, line~348). The longitudinal is an
in-plane displacement of brane material. Separately, the corpus has no
mechanistic account of what a gravitational wave is.

Second, a second cone is a departure only if matter couples to it. The
transverse wave equation is Lorentz-invariant with invariant speed
\(c_\gamma\), and matter in this model is built from those modes: a throat is
held open by a trapped brane-shear standing wave. This is not circular; the
wave equation supplies the invariance, so a standing wave built from it must
restructure under boost, with its cost diverging as \(v\to c_\gamma\). The
walk's statement that a second cone by itself means observers see no single
invariant speed was wrong.

Third, uniform flow is unobservable while gradients are observable. If uniform
flow were observable, absolute motion would be detected; if gradients were
not, there would be no gravity. ``The bulk flow is invisible because we are
made of it'' and ``density gradients bend light'' are statements about
different derivatives.

Fourth, light gravitates but does not dissipate, and the second fact follows
from S9. Light carries stress-energy as a co-moving disturbance without loss,
while photons from distant galaxies show that light must not leak. A photon
can lose energy only into an available mode, and the bulk has no transverse
mode. Photon stability over cosmological distance is therefore a consequence
of bulk shear-freeness. The genuine loss channel is mode conversion where the
medium is inhomogeneous, near matter and not in vacuum.

% NOT \stagefield{Verification}: that macro suppresses Verification in the
% default build.  The limitations and the meaning of PASS must stay visible.
\paragraph{Verification.}
Two independently written engines agree on every computed value. The blind
Mathematica audit
(\StageFile{research/pde\_ledger\_v3/mathematica/S11\_stray\_longitudinal\_mathematica\_audit.wl})
was written first, imports nothing, and reads no registry. The SymPy audit
(\StageFile{research/pde\_ledger\_v3/scripts/S11\_stray\_longitudinal\_sympy\_audit.py})
was written while the Mathematica file was quarantined, then compared with its
byte-identical committed blob \texttt{55620a7f}; the builder's git use was
checked as \texttt{status} and \texttt{diff} only.

The engines initially disagreed on one valuable point. Mathematica asserted a
conclusion about transverse coupling while SymPy refused to conclude it; an
independent review had already found the Mathematica line unconditional and
unearned. Mathematica was repaired to stop making the claim, not steered to
match SymPy. It then diagnosed its own circularity and named the missing input.
There were eight review legs: two for build directives, two per engine, and two
per repair.

A demonstrated control hole was also closed. Replacing \(R5\) by
\(c_L-\sqrt{\mu_R/\rho_{\mathrm{br}}}\), the exact wrong relation this step
exists to distinguish, had left all five gates green. A new assertion now
substitutes each derived root into the registry residual that records it; four
distinct corruptions were ablation-checked, and \(R4\) is covered as well
(\texttt{DEFECT\_REGISTER\#f-r5}).

The reported gates are acceptance \StatusText{match}
\((12\mathbin{\to}7,\ 12\mathbin{\to}7,\ 12\mathbin{\to}6)\);
dimensional homogeneity \StatusText{pass}, including homogeneous \(R5\) with
dimension \((1,-1,0)\) and zero undetermined; able-to-fail \StatusText{pass};
and 11 tests. Both engines print \texttt{VERDICT: PASS}. \textbf{Those strings
are verdicts only on the engines' internal checks; they are not verdicts on the
physics, the absent interface coupling, or the mode's physical status.}

\paragraph{Known limits --- recorded, not fixed.}
\begin{enumerate}[leftmargin=2.2em]
\item Two of the six dimensional checks in the Mathematica file are
definitional identities, \((X-2Y)+2Y=X\), and cannot fail. Only the
\(\rho_{\mathrm{br}}\) check obtained an independent route. The guard as a
whole is able to fail and four independent ablations were caught, but those
two entries are \textbf{not counted as coverage}.

\item \texttt{ASSERTION\_12} in the SymPy file defines
\(c:=\sqrt{X}\) and then asserts \(c^2-X=0\). It is identically zero by
construction, cannot fail, and is \textbf{not counted as coverage}; the new
registry control supersedes it in substance.

\item The Mathematica file disclaims scope on the transverse tag but not the
parallel tag, although both arise from the same substitution with the same
absent operator. The scope limit applies to both channels.

\item The SymPy script was repaired while untracked, so there is no committed
pre-repair baseline. A reviewer's \texttt{/tmp} copy and a full hand
re-derivation of every load-bearing value substituted for one. This violated
the rule to commit before anything destructive.

\item A review leg's search for \texttt{K\_br} incidentally returned two lines
of \StageFile{research/pde\_ledger\_v3/V3\_STEP\_PLAN.md}, which was on its
do-not-read list. The file was not opened, and all reported results had already
been derived.
\end{enumerate}


\section{S11b: The Linear Brane--Bulk Interface Coupling Law}
\label{step:S11b}

\stagefield{Part / anchor}{Part I --- Light. Sector~1, step~4.}
\claimstatus{\StatusText{Computed Interface Law / Controlled Stability Slice}}
\stagefield{Purpose}{Build the one object left open by S11: the linear law
that couples the brane's in-plane and thickness degrees of freedom to the two
bulk faces. The surviving result consists of the face response, the permeable
closure, the assembled balance laws, the admissibility and reciprocity
conditions, the uniform transverse null, and the breathing-mode stability
boundary.}

\paragraph{Status.}
The interface law is assembled on a uniform background. Its transverse block
is fixed there, and its breathing roots are fixed on the stated soluble slice.
The general non-uniform problem and the general dispersion remain open.

\paragraph{The object.}
S11 reduced three open questions --- whether the longitudinal mode radiates or
stays bound, whether light's confinement is unconditional, and whether the
second cone is observable --- to one missing law. This step supplies that law
by combining the bulk response to moving faces, interfacial transfer, the
brane's quadratic energy, and the slab balance equations in one linear system.

\paragraph{The bulk response to moving faces.}
For a disturbance with in-plane wavenumber \(k\) and frequency \(\omega\),
define
\[
  q^2=\frac{\omega^2}{c_{s0}^2}-k^2,
  \qquad
  Z=\frac{\omega\rho_m}{q_{\mathrm{out}}},
\]
where \(Z\) is pressure at a face divided by that face's outward normal
velocity. The outgoing branch gives
\begin{center}
\begin{tabular}{@{}L{0.20\textwidth}L{0.18\textwidth}L{0.50\textwidth}@{}}
\toprule
regime & response & meaning \\
\midrule
\(q^2>0\) & \(q_{\mathrm{out}}=q\),
  \(Z=\omega\rho_m/q\) & real radiation resistance \\
\(q^2<0\) & \(q_{\mathrm{out}}=i\alpha\),
  \(Z=-i\omega\rho_m/\alpha\) & reactive inertial loading \\
\(q^2=0\) & \(q_{\mathrm{out}}=0\), \(Z\to\infty\) & grazing
  singularity: a driven face has no solution, while an undriven face has a
  free amplitude \\
\bottomrule
\end{tabular}
\end{center}

\noindent
The two exterior half-spaces decouple, so the per-face response is the same
for the thickness and centre-shift combinations. Although the evanescent face
impedance carries the scale \(\rho_m/\alpha\), the assembled thickness
coordinate receives
\[
  m_{\mathrm{add}}^{(W,\mathrm{face})}=\frac{\rho_m}{2\alpha}
\]
from each face.

\paragraph{Permeable closure and its affinity.}
The assembled face law uses the affinity and velocity channels
\[
  J_\pm=\Lambda_A(\omega)\,\mathcal A_\pm
       +\Lambda_V(\omega)\,V_\pm,
  \qquad
  f_{X,\pm}=\Lambda_X(\omega)\,\mathcal A_\pm,
  \qquad
  \Lambda_I(\omega)=\frac{\Lambda_I^0}{1-i\omega\tau_I},
  \quad I\in\{A,V,X\}.
\]
Here \(\Lambda_A^0\), \(\Lambda_V^0\), and \(\Lambda_X^0\) are,
respectively, the DC (\(\omega\to0\)) affinity-channel flux,
velocity-channel flux, and reciprocal-traction-channel coefficients. The
thermodynamic driving variable is
\[
  \mathcal A_\pm=\mu_s-\frac{\delta p_\pm}{\rho_m}.
\]
On the slab-side-off slice \(\mu_s=0\), the affinity channel becomes a raw
pressure channel with the static map
\[
  \Lambda_p^0=-\frac{\Lambda_A^0}{\rho_m}.
\]
Its pressure contribution is therefore normalised by \(\rho_m\); replacing
\(\delta p/\rho_m\) by an unnormalised pressure is a different law. In the
equal-relaxation-time face-response slice the resulting pressure--velocity
parameterisation, with \(\tau_A=\tau_V=\tau\), gives
\[
  Z_{\mathrm{perm}}
  =\frac{\rho_m r+\Lambda_V^0}{y r-\Lambda_p^0},
  \qquad
  r=1-i\omega\tau,
  \qquad
  y=\frac{q_{\mathrm{out}}}{\omega}.
\]

\paragraph{Assembly by virtual displacement and balance law.}
The virtual-displacement rule fixes the face tractions, including their sign,
factor, and face count. The slab rows then follow from the interfacial mass
balance and from projecting four-dimensional continuity through a window
\(\Omega(w)\):
\[
  \partial_t\!\int\!\Omega\rho\,dw
  +\nabla_x\!\cdot\!\int\!\Omega j_x\,dw
  =\Omega(w_1)j^w(w_1)-\Omega(w_2)j^w(w_2)
   +\int j^w\Omega'\,dw.
\]
For an even window on the symmetric interval \([-L,L]\), even \(j^w\)
gives zero source exactly, whereas odd \(j^w\) gives
\(-2\Omega(L)j(L)+2\int_0^L j\Omega'\,du\), which need not vanish. With
outward relative flux
\(J_\pm=\rho_m(v_w-\partial_t\zeta_\pm)(\pm1)\), the symmetric combination is
net accretion \(-\left(J_++J_-\right)\), and the antisymmetric combination is
through-flow \((J_+-J_-)/2\). Two outgoing bulk amplitudes survive, and the
two face conditions fix them.

The resulting system couples the brane's in-plane displacement, its thickness
degree of freedom, interfacial transfer, and both bulk face amplitudes. Closing
its brane energy by symmetry yields \textbf{ten} independent quadratic
invariants, not the five supplied to the construction. The count is ten under
the shared specification's \S5 equivalence modulo total in-plane divergences;
the basis representative is not unique, and the two engines retain equivalent
representatives using, respectively, the traceless-symmetric strain square
and \((\nabla\!\cdot u)^2\) while spanning the same ten-dimensional quotient.
One engine initially over-counted eleven by omitting this quotient and was
corrected. In the \((\nabla\!\cdot u)^2\) representative retained here, the
five additional terms are
\[
  \frac{K_L}{2}(\nabla\!\cdot u)^2,
  \qquad
  \frac{A_\theta}{2}|\nabla\theta|^2,
  \qquad
  A_{\theta e}\,\nabla\theta\!\cdot\!\nabla e_W,
  \qquad
  D_\theta\theta(\nabla\!\cdot u),
  \qquad
  D_e e_W(\nabla\!\cdot u).
\]
Thus the five new coefficients are
\(K_L,A_\theta,A_{\theta e},D_\theta,D_e\). They enter the assembly through
\(m_\theta=B_\rho^{(3)}+A_\theta k^2\),
\(m_{\mathrm{cross}}=CW_0+A_{\theta e}k^2\), and the corresponding \(\pi\)
combinations.

\paragraph{Admissibility and reciprocity --- a classified region, not a gate.}
Requiring the interfacial power form to be non-negative, i.e.\ requiring the
interface to be \emph{passive}, gives
\[
  \Lambda_A^0\geq0,
  \qquad
  \left[
  \Lambda_V^0=\Lambda_X^0=0
  \quad\hbox{or}\quad
  \left(\Lambda_X^0=-\Lambda_V^0,\ \tau_V=\tau_X=0\right)
  \right],
\]
while \(\tau_A\) is unrestricted by the passivity conditions: there is no
relation beyond the supplied \(\tau_A\geq0\). A separate and
\emph{conditional} test --- Onsager--Casimir reciprocity, which holds only if
the underlying microdynamics is time-reversible --- fixes the full kernels,
\[
  \Lambda_X(\omega)=-\Lambda_V(\omega),
\]
reached by two independent partial-inversion routes that agree. Within the
passive region a nonzero velocity/reciprocal pair is therefore instantaneous.

\textbf{Both regions are classifications of parameter space, not acceptance
criteria, and neither is imposed anywhere in this ledger.} Two limits are
load-bearing. First, time-reversibility of the medium's microdynamics is
\emph{postulated nowhere in this model}, so the reciprocity relation is
conditional on an unstated assumption about a substructure the ledger does not
specify. Second, passivity and Onsager reciprocity are both properties of
fluctuations about a \emph{thermodynamic equilibrium}; this model's reference
state carries a steady background normal flow \(v_0\), making it a driven
steady state with a reservoir attached. About such a state neither property is
guaranteed, and non-reciprocal couplings of exactly the excluded form are
realised in driven laboratory media. Consequently the region boundary above
separates \emph{passive} from \emph{active} constitutive behaviour; it does not
separate possible from impossible.

\paragraph{Uniform transverse null.}
On a uniform background the transverse mode is exactly decoupled: in-plane
parity admits no \(e_W\leftrightarrow u_T\) bilinear. This coupling vanishes
identically, and the mode is non-dissipative unconditionally. In the
\((\nabla\!\cdot u)^2\) energy representative retained on this card, its
oscillator is
\[
  \rho_{\mathrm{br}}^0\omega^2=\mu_\perp k^2,
  \qquad \mu_\perp=\mu_R.
\]
It has a stable real frequency, \(\operatorname{Im}\omega=0\), only where
\(\mu_\perp=\mu_R\geq0\). No positivity is assumed for the moduli, so
\(\mu_\perp<0\) is admissible and gives the growing/decaying pair
\(\omega=\pm i k\sqrt{|\mu_\perp|/\rho_{\mathrm{br}}^0}\). (In the
cross-engine coefficient-basis comparison, the same physical stiffness is
\(\mu_R\) in the blind-Wolfram/card \((\nabla\!\cdot u)^2\) representative
and \(\mu_R+\mu_S/2\) in the SymPy traceless-symmetric-strain representative;
this is not the card's working formula.)
The dimension of the vanishing coupling is undetermined because a coefficient
that vanishes identically has no dimension. \textbf{This uniform transverse
null does not settle whether confinement is unconditional;} that is a
non-uniform question deferred to C.

\paragraph{Breathing-mode stability boundary.}
On the soluble slice \(k=0\), with impermeable faces and no reciprocal
coupling, the assembled determinant reduces to
\[
  \mu_W\omega^2+iR\omega-\frac{K_0}{W_0^2}=0,
  \qquad
  R=\frac{\rho_m c_{s0}}{2},
  \qquad
  K_0=B_\rho^{(3)}-2CW_0+k_WW_0^2.
\]
The mode is stable iff \(K_0>0\): \(K_0<0\) gives one growing and one decaying
root, \(K_0=0\) gives a static root, and \(K_0>0\) gives decay only.
Equivalently, growth occurs iff
\[
  C>\frac{B_\rho^{(3)}+k_WW_0^2}{2W_0}.
\]
The growing root is \textbf{not} an energy-conservation violation. The stored
energy has no minimum in that constrained direction, so the slab lowers its
energy by running away; the accounting closes, but the energy has no floor.
For \(K_0>0\), every nonzero imaginary part on this slice comes from
propagating bulk radiation resistance, which survives impermeability.

\paragraph{Departure --- the finite-memory velocity channel, and what it costs.}
On the pure-velocity branch \(\Lambda_p^0=0\), so the raw-pressure DC channel
is off and the coupling is velocity-driven. The face-response algebra admits
finite memory on this branch, with evanescent \(\operatorname{Re}Z\neq0\) iff
\(\omega\tau\neq0\), approaching the impermeable response as
\(\omega\tau\to\infty\). Its frequency profile is a Debye loss peak at
\(\omega\tau\sim1\), vanishing at both ends. On that branch the coupling
dissipates only at nonzero \(\omega\tau\), so slow disturbances see a leak
which fast disturbances do not.

This branch lies \emph{outside} the passive region: it is available only if the
interface draws on an energy reservoir. \textbf{That is a statement of price,
not of impossibility.} The model supplies a candidate reservoir --- the steady
background drain \(v_0\) --- so the branch is a live constitutive option whose
admission requires naming that reservoir and stating its power budget. It is
recorded here as a departure carrying that debt, and it is not asserted as a
prediction. The test that decides it is observational: the loss peak sits at
\(\omega\sim1/\tau\), and any \(\tau\) placing it in a band where anomalous
mechanical dissipation is measured and bounded is excluded by that measurement.
Converting this into a numerical bound requires the matter--brane coupling,
which this sector does not carry.

Two further preregistered predictions were refuted. The scalar-sector count is
two surviving internal degrees of freedom, not exactly one. The added inertia
on the thickness coordinate is \(\rho_m/(2\alpha)\) per face, not the predicted
\(\rho_m/\alpha\) per face.

\paragraph{Inputs and regime.}
The source of record for this card is the unified step record
\StageFile{research/pde\_ledger\_v3/steps/S11b\_interface\_coupling\_law.md}.
The historical
\StageFile{research/pde\_ledger\_v3/steps/S11bA\_interface\_response.md} and
\StageFile{research/pde\_ledger\_v3/steps/S11bB\_interface\_assembly.md} are
archival pre-unification history, not current sources. S11b is one unified
export-chain step subsuming A and B; C, the non-uniform transverse coupling,
is a separate later step.
The interface assembly and transverse result use a uniform background. The
breathing root statement is confined to the \(k=0\), impermeable,
no-reciprocal-coupling slice; it is not the general dispersion.

% NOT \stagefield{Verification}: that field is suppressed in the default
% build.  The scope of the checks and every limitation below must remain
% visible.
\paragraph{Verification scope.}
The SymPy engine imports the S11 \texttt{LEDGER}; the blind Wolfram engine
imports nothing and re-derives the objects, making it the only cross-engine
control. Their transcripts were compared by the frozen T7 comparator
\StageFile{research/pde\_ledger\_v3/scripts/S11b\_cross\_engine\_comparator.py},
with the recorded run in
\StageFile{research/pde\_ledger\_v3/scripts/out/S11b\_cross\_engine\_comparison.out}.
No compared object is a physics contradiction. The recorded disagreements are
emission format (SymPy \texttt{Tuple} versus Wolfram \texttt{Association}),
the coefficient-basis representative split, naming, a \(W_0\)-normalisation
convention, and a denominator-clearing form/domain artifact: the
\texttt{DEGENERATE\_LOCI} loci coincide on the regular domain and differ only
where the cleared denominator vanishes, including the grazing locus \(q=0\).
A few remaining cases are coverage gaps in which one engine solves an object
the other leaves open, not conflicts.

\paragraph{Known limits --- recorded, not fixed.}
\begin{enumerate}[leftmargin=2.2em]
\item \textbf{The background-flow correction is uncarried.} The rest-frame
bulk operator omits the convective correction
\(O(v_0|q_n|/\omega)\). The engines emit the relative term
\(-2qv_0/\omega\) (\texttt{RELATIVE\_TERM}), so first order fails where
\(|qv_0/\omega|\gtrsim1\). In the regime \(kc_{s0}\gg|\omega|\),
\(|q|\simeq k\), and failure requires
\[
  \frac{|v_0|}{c_{s0}}\frac{kc_{s0}}{|\omega|}\gtrsim1.
\]
Thus large \(kc_{s0}/|\omega|\) is necessary, not sufficient. The same steady
drain \(v_0\) is the candidate reservoir named for the departure; this is a
standing scope limit, not an active term in the assembled operator.

\item \textbf{The wall width is frozen.} The slab inertia
\(\rho_{\mathrm{br}}^0\equiv\rho_{4D}^0 W_0\) \emph{is} the imported
\(\rho_{\mathrm{br}}\); S11's wall width is held fixed, so \(\rho_{\mathrm{br}}\)
already is the slab-integrated three-inertia. The thickness response is taken at
this frozen \(W_0\). This is a recorded freeze, not a fix; a dynamical wall width
would carry terms this step does not.

\item \textbf{The supplied objects cannot be falsified by this build.} The
affinity, branch prescription, mass balance, virtual-displacement rule, and row
structure were supplied, so both engines transcribed them and an error would
land identically in both. Their verification comes from independent reviewer
derivations, \textbf{not} from the engines agreeing. This limitation is
load-bearing.

\item The breathing-root result is a \textbf{slice}, not the general
dispersion. There is no closed-form \(\omega(k)\), and no parameter-independent
root count exists.

\item The two-port balance check is blind to the face response, the closure,
and the affinity's \(\rho_m\) normalisation because they cancel between its two
sides. It tests the traction's sign, factor, and face count, not its analytic
structure; a transposed reciprocal kernel used identically on both sides
passes.

\item The \(Z_{\mathrm{perm}}\) acceptance value lies in the
equal-relaxation-time slice. The general \(\tau_A\neq\tau_V\) face response has
no independent standard.

\item Causality coverage is bounded. The pole inventory catches only
uncancelled, undisplaced advanced poles; cancelled and feedback-displaced poles
lie outside it, and zero-time and absent channels are indistinguishable from
passes.

\item The radiated-energy direction is inherited from the supplied
continuation. If that continuation is wrong, both engines are wrong together
and every decay/growth classification flips.

\item \StageFile{relations.yaml} remains empty. The admissibility condition and
reciprocity relation are emitted as review candidates, not inserted
automatically, because a text matcher does not verify that two notations are
equivalent.
\end{enumerate}


\section{S11c: Light at a Nonuniformity --- Partial}
\label{step:S11c}
\claimstatus{PARTIAL: no accepted nonuniform light-loss number}

S11c asked how much light escapes into other motion when the brane's
thickness or material properties change. Both reflected and transmitted
transverse light count as survival. The equations contain conversion
channels, but this campaign did not obtain a reliable loss fraction at a
defect. The canonical account is
\StageFile{steps/S11c_d_profile_conditioned_scattering.md}; the short handoff is
\StageFile{steps/S11c_PARTIAL_CLOSEOUT.md}.

\paragraph{What is established, conditionally.}
On the uniform background, light decouples from the bulk in the stated
linear model. Independent SymPy and Wolfram constructions support S11b's
result, with the coefficient identification and review limits recorded in
\StageFile{steps/S11b_interface_coupling_law.md} (transverse mode, comparison
and operational note). Stable real-frequency propagation requires
nonnegative transverse stiffness; a uniform result does not establish
confinement around a defect.

A later SymPy check found two selected light polarizations finite, carrying
current and not driving the bulk at the tested light/sound speed matches.
It fixes the bulk at rest, frequency \(\omega=3\), tangential momenta
\((1/5,1/10)\), held profiles and the bulk-density convention. The exact
matching sound speeds are \(\sqrt{3/2}\) on the left and
\(\sqrt{150/101}\) on the right. Claude and Grok cleared the selected method;
there was no fresh independent review of the executed result. The finite
samples and selected limits do not establish a complete mode census or
interval-wide regularity. See the coverage and limits in
\StageFile{_measurements/S11c_d_near_unity_uniform_continue_result.md}.

\paragraph{What is unresolved.}
The held-profile \(\omega=3\) scattering benchmark used rest bulk, the same
held-profile/density conventions and tangents, and a selected phase-speed
ratio approximately \(0.12\); the bare S9 coefficient ratio is separately
\(0.1\). Ten finite solves completed. In the final enlarged-domain,
regulator-halved setting, the two small signed current deficits were
\(-0.608218\) and \(+0.554792\) ppm. They lie within the empirical
\(1\) ppm numerical floor. That floor is \emph{not} a physical upper bound,
and the apparent gain is not evidence of a physical energy source.
Refinement/domain sensitivity and unfinished checks leave the result
\textbf{UNRESOLVED}. Both replacement reviewers returned
\textbf{NEEDS REVISION}; no matched-speed defect result follows. See the
signed-current table and controls in
\StageFile{_measurements/S11c_d_numerical_radiating_balance_report.md}.

\paragraph{Equation and interpretation limits.}
Four repairs to inertia, work orientation, pressure trace and thickness
coordinates have scoped review support, not full nonuniform-composition
clearance. The original b/c2 verification describes pre-repair equations;
their records' ``Changed after close'' sections carry the later limits.
The c2 fold sets the direct three-leg \([0,2]\) entry to zero while keeping
the iterated contribution. Retaining \(\eta\sigma_W\) means order counting
does not justify that zero. Whether the complete closed response needs a
nonzero direct term remains \textbf{UNRESOLVED}; no benchmark correction
has been established. See ``The remaining mixed-term issue'' in
\StageFile{_measurements/S11c_upstream_repair_review_joint_disposition.md}.

The clean-condition packet measured a symmetry selection rule on SymPy
only. Its no-leak interpretation requires the drain frozen; it does not
transfer to live conversion. The broader packet remained uncleared
(Claude: ``not clear.''; Grok: ``not cleared.''). See round~5 of
\StageFile{directives/_measurements/S11c_d_clean_condition_review_disposition.md}.
Uniform confinement is derived within the supplied model; nonuniform
confinement and the material admissibility of its rotational stiffness
remain open. Intermediate source, packet, current and receiving checks do
not replace a complete physical work and outgoing-power balance.

\paragraph{Handoff.}
S12 is the next planned subject: dynamical bulk-to-brane order conversion
and its separate boundary data. It does not require a numerical S11c loss
factor. Equation/computation debts stay with S11c-b/c2/d if reopened;
S11c-e owns the conversion observable and strong-edge interpretation.
Q2/Q3/S22 owns real throat support and charge response; S22/R10 owns the
proposed speed-ratio derivation, with calibration retained in S20a.
S22's conversion/birefringence tests need nonuniform response and a supplied
profile, not a loss number presumed from this partial close.
These ownership boundaries are in \StageFile{V3_STEP_PLAN.md} (S12, charge
preamble, S20a and S22). The throat/EM notes indexed in
\StageFile{docs/light_em_investigation_handoff.md} at the repository root
remain exploratory, paused inputs to S22/Q2, not adopted material laws.
